\section{Minecraft as a Construction Substrate}
\label{sec:minecraft-mechanics}

Minecraft's world is a three-dimensional grid of one-metre cubes, or \emph{blocks}; a block
occupies exactly one grid cell and, depending on its type, may carry a small amount of state
directly within that cell (its \emph{block state}), such as which of the six axis-aligned
directions it is \emph{facing} at the moment of placement. Every two-terminal component in
this mod reuses exactly this mechanic to determine which two of its six faces serve as its
electrical leads (Section~\ref{sec:architecture}): a resistor placed while facing north has
its leads on its north and south faces, the same behavior the base game's own oriented blocks
(a piston, a dispenser) already exhibit, so no new placement mechanic had to be introduced.

The base game already provides an entire circuit-like subsystem, redstone: an integer signal
strength from 0 to 15 that decreases by one per block of dust traversed and updates through a
discrete neighbor-notification event fired once per tick, with no built-in notion of voltage,
current, capacitance, or continuous time. Despite this limitation, redstone is widely used,
both informally and as a subject of academic interest~\cite{nebel2016mining}, as an entry
point into digital logic and rudimentary computer architecture, and players commonly describe
redstone circuits in explicitly electrical language despite the underlying mechanic being
discrete and boolean~\cite{dezuanni2015redstone}. This motivates a specific design decision:
Minecraft's building metaphor already appears to prime players to reason electrically, yet
redstone cannot itself serve as the vehicle for teaching \emph{analog} behavior, since it has
no continuous quantity to teach with -- community ``analog redstone'' techniques exist, but
only approximate a continuous value as a coarse, sixteen-level discretization.

Minecraft's modding interfaces expose a richer construct than a plain block's fixed state,
the \emph{block entity}: an object holding arbitrary mutable data that, if requested, executes
code once per game tick (20~Hz, the game's own simulation rate). Every component, wire node,
ground reference, and configuration module in this mod is a block entity for precisely this
reason, locking the mod's simulation cadence to the game's own tick rate without a separate
timing mechanism. The mod's own contribution is deliberately narrow: a second, parallel
notion of adjacency (whether a face is electrically conductive toward its neighbor,
Section~\ref{sec:architecture}) alongside the game's existing physical adjacency, supplying
the resulting graph to a genuine solver every tick rather than encoding behavior directly
into the blocks. Redstone is reused only for the one concept it is genuinely suited to
represent here: a boolean on/off switch gating whether a voltage source is connected at all
(Section~\ref{sec:pedagogy}), not as an analog signal it was never designed to carry.
